\section{Reproduction package}
\label{app:repro}

Every number, table entry and curve in this paper is generated by the scripts
listed below from the two material files listed below.  No value is transcribed
by hand, and no plotted quantity is computed inside the plotting script.

\subsection{Contents of the package}
\label{app:repro_contents}

\begin{table}[h]
\centering
\caption{Files required to reproduce the paper.  Paths are relative to the
archive root and are broken across lines by \texttt{url} conventions, so no
entry can run past the right margin.}
\label{tab:package}
\small
\begin{tabularx}{\textwidth}{@{}>{\raggedright\arraybackslash\ttfamily\footnotesize}p{0.34\textwidth}>{\raggedright\arraybackslash}X@{}}
\toprule
\normalfont\itshape Path & \normalfont\itshape Role \\
\midrule
\path{a0/materials.json} & material constants and their provenance; the model cannot start without it \\
\path{paper_a_study/model.py} & spectral plate solver, spatial and temporal routes, edge matching, power and dissipation \\
\path{paper_a_study/build_figure_data.py} & solves the physics and writes every CSV/JSON used by the figures \\
\path{paper_a_study/certify_peak.py} & degree-refinement study of the two reported values (\cref{tab:refinement}) \\
\path{paper_a_study/reclassify_csv.py} & re-derives the admissibility flags of the figure CSVs without re-solving \\
\path{paper_a_study/generate_final_figures.py} & reads the CSVs and renders the ten vector PDFs \\
\path{paper_a_study/figure_style.py} & shared typography and colour map \\
\path{paper_a_study/make_atlas.py} & two-page contact sheet of the ten figures \\
\path{paper_a_study/test_model.py} & the 20 automated checks of \cref{sec:numerics_verification} \\
\path{paper_a_study/verify_dimensions.py} & dimensional and constitutive audit, including the strain-gradient contraction test \\
\path{paper_a_study/quasistatic_check.py} & the ratios of \cref{sec:governing_electric} \\
\path{paper_a_study/flexo_control.py} & the measured flexoelectric contribution of \cref{sec:results_scope} \\
\path{paper_a_study/phase_speed_audit.py} & phase-speed and wavenumber reconciliation of \cref{sec:governing_electric}; exits 0 when the text agrees with the solver and non-zero when it does not \\
\path{paper_a_study/modal_flux_audit.py} & eigenvector-rescaling and degree-refinement proof of $R_B$ (\cref{sec:results_peak}) \\
\path{paper_a_study/mechanism_controls.py} & the matched controls of \cref{tab:controls} and the branch-consistent $H$ sweep \\
\path{paper_a_study/corner_assessment.py} & corner-band energy share and incident-wavenumber regression behind \cref{sec:disc_limitations} \\
\path{paper_a_study/run_all_checks.py} & the whole battery, one log \\
\path{paper_a_study/figure_data/*.csv} & the twelve data files behind the figures \\
\path{paper_a_study/figure_data/*.json} & panel-source map, figure manifest, field metrics, refinement study, controls and corner assessment \\
\path{main.tex, sections/, figures/} & this manuscript and the ten vector PDFs \\
\bottomrule
\end{tabularx}
\end{table}

The archive deliberately does not contain \texttt{.npy} or \texttt{.npz} files.
Two legacy scripts that are superseded by
\texttt{build\_figure\_data.py}/\path{generate_final_figures.py} used such
caches; they are not part of the reproduction path and are not shipped.  Every
indispensable input is a JSON, a CSV or a Python source file, so a clean
extraction reproduces the results without any binary intermediate.

\subsection{Commands}
\label{app:repro_commands}

From a clean extraction of the archive, with a Python 3.11 environment
containing \texttt{numpy}, \texttt{scipy} and \texttt{matplotlib}:

\begin{verbatim}
# exit status 0 = every check passed
python paper_a_study/run_all_checks.py            # battery, ~25 s
python -m unittest paper_a_study.test_model -v    # 20 tests, ~9 s
python paper_a_study/build_figure_data.py         # CSV/JSON, ~6 min
python paper_a_study/reclassify_csv.py            # flags, ~1 s
python paper_a_study/certify_peak.py              # refinement, ~4 min
python paper_a_study/generate_final_figures.py    # ten PDFs, ~20 s
python paper_a_study/make_atlas.py                # contact sheet
python paper_a_study/modal_flux_audit.py          # rescaling proof
python paper_a_study/mechanism_controls.py        # controls, ~1 min
python paper_a_study/corner_assessment.py         # corner check, ~20 s
python paper_a_study/phase_speed_audit.py         # 0 = text agrees
\end{verbatim}

The order matters only in that \path{generate_final_figures.py} needs the
CSV files.  \path{certify_peak.py} and \path{reclassify_csv.py} are
independent of each other and of the plotting step.  The scripts resolve every
path relative to their own location (the value of \path{__file__}) and to the
repository root, so they can be run from any working directory.

Exit statuses are part of the interface and are stated here because they are
easy to mis-state.  \path{run_all_checks.py}, \path{reclassify_csv.py},
\path{certify_peak.py}, \path{generate_final_figures.py},
\path{modal_flux_audit.py}, \path{mechanism_controls.py} and
\path{corner_assessment.py} return 0 when their own checks pass and non-zero
when one fails.  \path{phase_speed_audit.py} returns 0 when every quantity
quoted in \cref{sec:governing_electric} agrees with the recomputation from the
material file, and non-zero when one does not; it is a consistency check
between the manuscript and the solver, so a non-zero status there is a defect
report, not a normal outcome.  None of the scripts is expected to fail on the
shipped tree: the last full run is recorded in
\path{validation/verification_log.txt}.

\subsection{Environment used for the reported runs}
\label{app:repro_env}

Python 3.11.2, NumPy 2.4.6, SciPy 1.17.1, Matplotlib 3.11.2.  The solver uses
only dense linear algebra (\texttt{numpy.linalg}, \texttt{scipy.linalg.svd} and
the generalised eigenproblem), so results are reproducible across BLAS
implementations to round-off.  The only stochastic element in the package is
the random state used inside two of the verification tests; no production
result depends on it.

\subsection{Provenance records}
\label{app:repro_provenance}

\texttt{a0/materials.json} records every constant twice: once as a value and
once as a provenance tag.  The tags distinguish literature values, values
derived through the Einstein relation, and synthetic parameters chosen for this
study, exactly as \cref{tab:provenance} reports.  The file also records the
values that were deliberately \emph{not} assigned, namely the flexoelectric
tensor of a specific specimen, the strain-gradient tensor of a specific
specimen, the $H$ tensor of any specimen, the sample doping, the sample
thickness and the surface properties.  These are the quantities for which no
defensible literature value was found; they are marked as such rather than
filled with a plausible number.

\path{figure_data/figure_manifest.json} maps each of the ten figures to its
data file and to the script that produced it, and
\path{figure_data/figure_panel_sources.json} maps each individual panel to
the columns of the corresponding CSV.  Together they allow any single panel to
be traced back to the solve that produced it.

\subsection{What is not reproducible here}
\label{app:repro_limits}

The visual appearance of the figures is reproducible only up to the Matplotlib
and font versions, because the layout is tuned to the reported versions.  The
manuscript itself was validated statically: environment and bracket balance,
cross-reference resolution, citation-key resolution and the presence of every
included figure file were checked by script, but no TeX engine was available at
the time of writing, so the PDF was not compiled in this environment.  A reader
who builds the manuscript must therefore treat the typeset output as
unverified and the content as validated.
